\documentclass[11pt,a4paper]{article}
\usepackage[utf8]{inputenc}
\usepackage[T1]{fontenc}
\usepackage{mathpazo}
\usepackage[scaled=0.92]{helvet}
\usepackage[english]{babel}
\usepackage{microtype}
\usepackage[a4paper,margin=2.2cm]{geometry}
\usepackage{xcolor}
\definecolor{anscaption}{HTML}{6B7280}
\usepackage[colorlinks=true,linkcolor=blue,urlcolor=blue,pdfborder={0 0 0}]{hyperref}
\usepackage{fancyhdr}
\pagestyle{fancy}
\fancyhf{}
\renewcommand{\headrulewidth}{0pt}
\fancyfoot[L]{\footnotesize\color{anscaption}CSE 341: ANS Retrospective (D5, Part 2) \(\cdot\) 220104004926 Yusuf Hamza Çelebi}
\fancyfoot[R]{\footnotesize\color{anscaption}Page \thepage}
\setlength{\parindent}{0pt}
\setlength{\parskip}{5pt}
\newcommand{\code}[1]{\texttt{#1}}
\begin{document}
\begin{center}
{\footnotesize\color{anscaption}CSE 341: Concepts of Programming Languages\par}
{\huge\bfseries ANS: Alchemist's Notation Script\par}
{\large Retrospective \& Self-Assessment (D5), Part 2\par}
{\footnotesize\color{anscaption}Student: 220104004926 Yusuf Hamza Çelebi}
\end{center}

\section*{Retrospective}
If I started again, I would design the Part 2 type system before writing the parser examples. In Part 1 we allowed programs that parsed but were chemically awkward, especially the first \code{compute\_yield} shape that tried to look up the same compound in both reactants and products. Part 2 forced that design to become precise: the final function separates reactant and product compounds.

The hardest decision was how much chemistry to put in the language instead of user functions. We kept formula parsing, compound records, reaction sides, and \code{check = balanced} in the language because those features make ANS a real DSL. We scoped out a larger unit algebra system; quantities are recognized and \code{mol(x)} extracts numeric amounts, but the interpreter does not do full dimensional analysis. That is a reasonable limit for this project.

\section*{Self-Assessment}
\textbf{Correctness: B.} The three example programs execute, malformed literals and invalid assignment targets are rejected, type errors are caught before execution, and required runtime errors are reported. Some edge cases such as complete unit algebra are intentionally limited.

\textbf{Design Justification: B.} The final D1 ties major choices to Sebesta terms and names trade-offs, especially strong typing, static scoping, and assignment as a statement.

\textbf{Use of Sebesta Framework: B.} The report uses EBNF, operational semantics, type-system terminology, binding time, scope, and lifetime accurately enough to defend in the exam.

\textbf{Originality: A.} ANS is clearly chemistry-specific. Formula literals, checked reactions, and reaction-side indexing are not just generic scripting features with chemistry names.

\textbf{Code Quality: B.} The lexer, parser, AST, semantic checker, and interpreter are separated. The runtime is larger than ideal, but the final tests now fail loudly when an archive is incomplete or a negative case is accepted.

\textbf{Exam Readiness: B.} I can explain the parser and the chemistry runtime decisions. I still need to rehearse the exact operational semantics and type-equivalence examples before the written exam.

\end{document}
